\chapter{Implied Volatility and Its Surface}\label{ch:dv:implied-volatility-and-its-surface}

Before October 1987 the implied volatilities of S\&P~500 index options were
roughly the same at every strike: a flat line, as the \omterm{def:dv:black-scholes-three-ways:model}{Black--Scholes model}
says they should be. After the crash, in which the index fell by a fifth within
two days, low strikes began to trade at higher volatilities than high strikes,
and they have done so ever since. On a quiet day now, a three-month put 10\%
below the market costs 20 volatility while the call 10\% above costs 14; a put
20\% below costs 27, and prices a fall that the at-the-money volatility of 16
would call a one-in-three-hundred event as a one-in-fifty event. The model has
not changed; the prices have. Since then the model has been used as a unit of
measure: every option is quoted by the volatility that the formula needs to
reproduce its price, and the table of those numbers across strikes and expiries
is what an options desk looks at, fits, trades and risk-manages. This chapter
describes that table, the rules it must obey, and what it says about the
distribution the market is pricing.

\section{The surface as the market's price list}

Implied volatility is defined in One Quant Book 1, chapter 25, with the skew
across strikes and the term structure across expiries. Put together they form a
surface, and two changes of coordinates make it comparable across expiries and
underlyings.

\begin{definition}[Volatility surface, log-moneyness, total implied variance]\label{def:dv:implied-volatility-and-its-surface:surface}
The \emph{volatility surface}\index{volatility surface} of an underlying is the
function $(K,T)\mapsto\sigma_{\mathrm{imp}}(K,T)$ that gives, for every strike and
expiry, the implied volatility of the European option, computed with Black's formula
on the forward $F_{0,T}$ of that expiry. Its natural coordinates are the
\emph{log-moneyness}\index{log-moneyness} $k=\ln(K/F_{0,T})$, which measures a strike
against the forward of its own expiry, and the \emph{total implied
variance}\index{total implied variance} $w(k,T)=\sigma_{\mathrm{imp}}^2T$, the variance
of the log-return to expiry that the price implies.
\end{definition}

In these coordinates Black's formula depends on $(k,w)$ only, up to the factor
$P(0,T)F_{0,T}$: $C/(PF)=\Phi(-k/\sqrt w+\sqrt w/2)-e^k\Phi(-k/\sqrt w-\sqrt w/2)$.
A surface is then a family of smiles $k\mapsto w(k,T)$, one per expiry, and the
static no-arbitrage conditions of the next sections become conditions on $w$.

\begin{omfigure}
\begin{tikzpicture}
\begin{axis}[width=10.5cm,height=5cm,
  xlabel={strike (spot 100)},ylabel={implied volatility (\%)},
  tick label style={font=\small},label style={font=\small},
  xmin=60,xmax=130,ymin=8,ymax=55,grid=major,grid style={gray!25},
  legend columns=4,legend style={font=\footnotesize,at={(0.5,-0.3)},anchor=north,draw=none,fill=none,/tikz/every even column/.append style={column sep=8pt}},
  table/col sep=comma]
\addplot[omProp,thick] table[x=strike,y=m1]{figdata/derivatives/07-implied-volatility-and-its-surface/smiles.csv};
\addplot[omDef,thick] table[x=strike,y=m3]{figdata/derivatives/07-implied-volatility-and-its-surface/smiles.csv};
\addplot[omThm,very thick] table[x=strike,y=y1]{figdata/derivatives/07-implied-volatility-and-its-surface/smiles.csv};
\addplot[gray,very thick,dashed] table[x=strike,y=y2]{figdata/derivatives/07-implied-volatility-and-its-surface/smiles.csv};
\legend{1 month,3 months,1 year,2 years}
\end{axis}
\end{tikzpicture}

\omcaption{The chapter's equity-index surface: four smiles of one surface. Short
expiries are steep and low at the money; long expiries are flatter and higher. The
90--110 skew is 10.8 volatility points at one month, 6.3 at three months, 3.3 at one
year and 2.5 at two. Illustrative parameters, in line with the shapes described in
the text.}\label{fig:dv:implied-volatility-and-its-surface:smiles}
\end{omfigure}

The surface is also the desk's price list. A quote for any European option on the
underlying, however unusual its strike or expiry, is read from it; a structured product
priced by a model calibrated to it inherits its skew; and the risk of the whole book is
measured against moves of it.

\section{Smile phenomenology by asset class}

The shapes differ by market, and the differences have reasons.

\begin{itemize}
\item \textbf{Equity indices}: a steep negative skew, steeper for short expiries and
decaying roughly like a power of the expiry (chapter 12), with the at-the-money level
low in calm markets and the term structure upward-sloping. Investors buy index puts as
insurance, crashes happen downward, and volatility rises when prices fall.
\item \textbf{Single stocks}: skewed the same way but less, with more curvature: studies
of individual equity options find their risk-neutral distributions far less negatively
skewed than the index's. The index skew adds the tendency of stocks to fall together,
the correlation skew of chapter 17.
\item \textbf{Currencies}: nearly symmetric smiles around the at-the-money level, tilted
by the risk reversal towards the currency that is feared to fall (One Quant Book 2,
chapter 19).
\item \textbf{Commodities}: often skewed upwards. A supply shock raises the price and the
volatility together, and most commodities show higher volatility after price rises, the
reverse of the equity pattern; crude oil is the documented exception
(\cref{fig:dv:implied-volatility-and-its-surface:classes}).
\item \textbf{Interest rates}: quoted in normal volatility (One Quant Book 2, chapter 13),
with smiles modelled by SABR (chapter 11) and a cube of expiries, tenors and strikes
built in One Quant Book 6, chapter 5.
\end{itemize}

\begin{omfigure}
\begin{tikzpicture}
\begin{axis}[width=10.5cm,height=5cm,
  xlabel={log-moneyness $k$},ylabel={implied volatility (\%)},
  tick label style={font=\small},label style={font=\small},
  xmin=-0.3,xmax=0.3,ymin=5,ymax=42,grid=major,grid style={gray!25},
  legend columns=3,legend style={font=\footnotesize,at={(0.5,-0.3)},anchor=north,draw=none,fill=none,/tikz/every even column/.append style={column sep=8pt}},
  table/col sep=comma]
\addplot[omThm,very thick] table[x=k,y=equity]{figdata/derivatives/07-implied-volatility-and-its-surface/asset_classes.csv};
\addplot[omDef,very thick] table[x=k,y=fx]{figdata/derivatives/07-implied-volatility-and-its-surface/asset_classes.csv};
\addplot[omProp,very thick,dashed] table[x=k,y=commodity]{figdata/derivatives/07-implied-volatility-and-its-surface/asset_classes.csv};
\legend{equity index,currency pair,commodity}
\end{axis}
\end{tikzpicture}

\omcaption{Three-month smiles by asset class, as functions of \omterm{def:dv:implied-volatility-and-its-surface:surface}{log-moneyness}: an equity
index skewed down, a currency pair nearly symmetric, a commodity skewed up.
Illustrative shapes.}\label{fig:dv:implied-volatility-and-its-surface:classes}
\end{omfigure}

\section{Sticky strike, sticky delta and what the desk assumes}

A surface is a snapshot. When the spot moves, something must be assumed about how the
surface moves with it, and the assumption changes the hedge.

\begin{definition}[Sticky strike and sticky delta]\label{def:dv:implied-volatility-and-its-surface:sticky}
Under \emph{sticky strike}\index{sticky strike} the implied volatility of each strike
and expiry stays fixed when the spot moves: the smile, drawn against strike, does not
move. Under \emph{sticky delta}\index{sticky delta} (also sticky moneyness) the implied
volatility of each moneyness $K/F$ stays fixed: the smile moves with the forward, and
the at-the-money volatility does not change.
\end{definition}

The two rules are regimes rather than laws: studies of index volatilities find periods
in which each describes the market better, and a third in which the at-the-money
volatility moves more than either predicts (Derman). They matter for three reasons.
They give different P\&L for the same move: after a 5\% fall of the index, the
three-month put struck at 100 is worth 6.00 under \omterm{def:dv:implied-volatility-and-its-surface:sticky}{sticky strike} and 5.78 under \omterm{def:dv:implied-volatility-and-its-surface:sticky}{sticky
delta}, from 3.08 before (\cref{fig:dv:implied-volatility-and-its-surface:regimes}).
They give different deltas.

\begin{proposition}[Smile-adjusted delta]\label{prop:dv:implied-volatility-and-its-surface:delta}
If the implied volatility of a fixed strike moves with the spot as $\partial\sigma_{\mathrm{imp}}/\partial S$,
the option's total delta is $\Delta_{\mathrm{BS}}+\mathcal V\,\partial\sigma_{\mathrm{imp}}/\partial S$.
Under \omterm{def:dv:implied-volatility-and-its-surface:sticky}{sticky strike} the correction is zero; under \omterm{def:dv:implied-volatility-and-its-surface:sticky}{sticky delta}, with the smile a function
of $k=\ln(K/F)$, $\partial\sigma_{\mathrm{imp}}/\partial S=-\sigma_{\mathrm{imp}}'(k)/S$, which is
positive for a negative skew.
\end{proposition}
\begin{proof}
Differentiate $V=C_{\mathrm{BS}}(S,\sigma_{\mathrm{imp}}(K,T;S))$ in $S$; with
$\sigma_{\mathrm{imp}}=\Sigma(\ln K-\ln F)$ and $F$ proportional to $S$,
$\partial\sigma_{\mathrm{imp}}/\partial S=-\Sigma'(k)/S$.
\end{proof}

For the three-month call struck at 100 on the chapter's surface, the Black--Scholes delta
is 0.533 and the sticky-delta delta 0.593: a desk that assumes one regime and lives in the
other is under- or over-hedged by six shares per hundred options. And they give different
exotic prices, because a model of the surface's dynamics (local volatility, chapter 9;
stochastic volatility, chapter 10) implies one regime or the other; chapter 11 measures
the difference.

\begin{omfigure}
\begin{tikzpicture}
\begin{axis}[width=10.5cm,height=5cm,
  xlabel={strike},ylabel={3-month implied volatility (\%)},
  tick label style={font=\small},label style={font=\small},
  xmin=80,xmax=120,ymin=10,ymax=28,grid=major,grid style={gray!25},
  legend columns=2,legend style={font=\footnotesize,at={(0.5,-0.3)},anchor=north,draw=none,fill=none,/tikz/every even column/.append style={column sep=8pt}},
  table/col sep=comma]
\addplot[omDef,very thick] table[x=strike,y=before]{figdata/derivatives/07-implied-volatility-and-its-surface/regimes.csv};
\addplot[omThm,very thick,dashed] table[x=strike,y=sticky_delta]{figdata/derivatives/07-implied-volatility-and-its-surface/regimes.csv};
\draw[omDef] (axis cs:100.37,16.36) circle[radius=2.5pt];
\draw[omThm] (axis cs:95.36,16.36) circle[radius=2.5pt];
\draw[omDef,dotted,thick] (axis cs:95.36,10) -- (axis cs:95.36,18.08);
\node[font=\footnotesize,anchor=west,omDef,align=left] at (axis cs:96,21.5) {new at-the-money strike 95.4:\\18.1 under sticky strike};
\legend{{before, and after the fall under sticky strike},after a 5\% fall under sticky delta}
\end{axis}
\end{tikzpicture}

\omcaption{The three-month smile before and after a 5\% fall of the index. Under \omterm{def:dv:implied-volatility-and-its-surface:sticky}{sticky
strike} the curve stays where it was and the new at-the-money strike (95.4) reads 18.1
volatility on it; under \omterm{def:dv:implied-volatility-and-its-surface:sticky}{sticky delta} the curve shifts left with the forward and the
at-the-money volatility stays at 16.4. Round markers: at-the-money points. Data: the
tutorial.}\label{fig:dv:implied-volatility-and-its-surface:regimes}
\end{omfigure}

\section{Static no-arbitrage conditions}

Not every table of numbers is a surface. Chapter 1 listed the model-free inequalities on
the prices of one expiry; across expiries one more family is needed.

\begin{definition}[Static arbitrage, butterfly and calendar arbitrage]\label{def:dv:implied-volatility-and-its-surface:static}
A \emph{static arbitrage}\index{static arbitrage} in a set of option quotes is a portfolio
of the options, the forward and bonds, bought once and held to the expiries, that costs
nothing and pays a non-negative amount in every state, positive in some. A
\emph{butterfly arbitrage}\index{butterfly arbitrage} is one inside an expiry: call prices
that are not convex in the strike, so that a butterfly has a negative price. A
\emph{calendar arbitrage}\index{calendar arbitrage} is one across expiries: a longer-dated
option cheaper than the shorter-dated one at the same forward-moneyness.
\end{definition}

The absence of call-spread, butterfly and \omterm{def:dv:implied-volatility-and-its-surface:static}{calendar arbitrage} is sufficient to exclude all
\omterm{def:dv:implied-volatility-and-its-surface:static}{static arbitrage} from a surface of quotes on one underlier (Carr and Madan). In total
variance the conditions are simple.

\begin{proposition}[Static arbitrage in total variance]\label{prop:dv:implied-volatility-and-its-surface:static}
With deterministic rates and dividends proportional to the share: (i) there is no \omterm{def:dv:implied-volatility-and-its-surface:static}{calendar
arbitrage} if and only if $w(k,T)$ is non-decreasing in $T$ for every $k$; (ii) a smile
$k\mapsto w(k)$ is free of \omterm{def:dv:implied-volatility-and-its-surface:static}{butterfly arbitrage} if and only if the call prices it implies
are convex in the strike, which holds when
\[
g(k)=\Bigl(1-\frac{kw'}{2w}\Bigr)^2-\frac{w'^2}{4}\Bigl(\frac1w+\frac14\Bigr)+\frac{w^{\prime\prime}}{2}\ \ge0
\]
for every $k$, together with call prices that tend to zero for large strikes.
\end{proposition}
\begin{proof}[Partial proof]
(i) At fixed $k$, the normalised call $C/(PF)$ is an increasing function of $w$ alone; a
calendar spread at fixed forward-moneyness is a call on a martingale at two dates, whose
value increases with the date by Jensen's inequality, hence $w$ must increase. (ii) The
second strike derivative of the call equals the density of the next section, and computing
it through $w(k)$ gives $g(k)$ times a positive factor $\varphi(d_2)/\sqrt w$; the full
statement is proved by Gatheral and Jacquier.
\end{proof}

A single bad quote is enough to break a surface. On the chapter's surface, raising the
three-month volatility at $k=-0.1$ by three points makes $g$ negative at six points of that
smile, down to $-2.06$: the butterfly centred there would have a negative price. Lowering
the six-month at-the-money quote by seven points makes $w$ fall between three and six months:
a \omterm{def:dv:implied-volatility-and-its-surface:static}{calendar arbitrage}. The surface builder of this chapter checks both before any model is
calibrated to the surface (chapter 8).

\section{The risk-neutral density}

\begin{definition}[Risk-neutral density and the Breeden--Litzenberger formula]\label{def:dv:implied-volatility-and-its-surface:density}
The \emph{risk-neutral density}\index{risk-neutral density} of $S_T$ is the density
$q_T(s)$ of the underlying at $T$ under the risk-neutral measure (One Quant Book 4,
chapter 5); it is the continuous version of the \omterm{def:dv:no-arbitrage-and-the-fundamental-theorems:ad}{state prices} of chapter 1, divided by the
discount factor. The \emph{Breeden--Litzenberger formula}\index{Breeden--Litzenberger formula}
reads it off call prices:
\[
q_T(K)=\frac{1}{P(0,T)}\,\frac{\partial^2C(K,T)}{\partial K^2}.
\]
\end{definition}
\begin{proof}[Derivation]
The call is the discounted expected payoff,
\[
C(K,T)=P(0,T)\int_K^\infty(s-K)\,q_T(s)\,ds .
\]
Differentiating once in $K$ gives $-P(0,T)\int_K^\infty q_T(s)\,ds$, minus the discounted
digital, and once more gives the formula.
\end{proof}

The first derivative is already a statement about the skew: the price of a digital call is
$-\partial C/\partial K=P(0,T)\Phi(d_2)-\mathcal V\,\partial\sigma_{\mathrm{imp}}/\partial K$,
the Black--Scholes digital plus a skew correction. With a negative skew the correction is
positive: the three-month at-the-money digital call on the chapter's surface is worth 0.558,
against 0.498 without the skew.

\begin{example}[The crash premium]\label{ex:dv:implied-volatility-and-its-surface:crash}
On the chapter's surface the three-month forward is 100.37 and the at-the-money volatility
16.36\%. With a flat surface at that level the risk-neutral probability that the index ends
below 80 is 0.31\%; with the skew, whose 80 strike is at 26.8 volatility, it is 1.89\%, six
times more (\cref{fig:dv:implied-volatility-and-its-surface:density}). The 80 put costs 0.006
on the flat surface and 0.219 on the skewed one.
\end{example}

\begin{omfigure}
\begin{tikzpicture}
\begin{axis}[width=10.5cm,height=5cm,
  xlabel={index level in 3 months},ylabel={risk-neutral density},
  tick label style={font=\small},label style={font=\small},
  xmin=55,xmax=140,ymin=0,ymax=0.058,grid=major,grid style={gray!25},
  yticklabel style={/pgf/number format/.cd,fixed,precision=2},scaled y ticks=false,
  legend columns=2,legend style={font=\footnotesize,at={(0.5,-0.3)},anchor=north,draw=none,fill=none,/tikz/every even column/.append style={column sep=8pt}},
  table/col sep=comma]
\addplot[omThm,very thick] table[x=s,y=skewed]{figdata/derivatives/07-implied-volatility-and-its-surface/density.csv};
\addplot[omDef,thick,dashed] table[x=s,y=flat]{figdata/derivatives/07-implied-volatility-and-its-surface/density.csv};
\draw[gray,dotted,thick] (axis cs:80,0) -- (axis cs:80,0.058);
\legend{skewed surface,flat at the at-the-money volatility}
\end{axis}
\end{tikzpicture}

\omcaption{Risk-neutral densities of the index in three months from Breeden--Litzenberger
applied to call prices: with the chapter's skew (solid) and with a flat surface at the same
at-the-money volatility (dashed). The skew moves probability into the left tail and makes
the density's peak sit above the forward; the probability below 80 (dotted) rises from
0.31\% to 1.89\%. Data: the tutorial.}\label{fig:dv:implied-volatility-and-its-surface:density}
\end{omfigure}

The density is a price, not a forecast. It includes what investors pay for protection, and
its left tail is fatter than the real-world frequency of crashes if insurance commands a
premium, as chapter 14 measures on variance. What it is exactly is the portfolio of options
that pays one unit if the index ends in a small interval around $s$: the \omterm{def:dv:no-arbitrage-and-the-fundamental-theorems:ad}{Arrow--Debreu
security} of chapter 1 made continuous, and replicable from the strip of listed options.

\section{Tutorial: building and checking a surface}\label{tut:dv:implied-volatility-and-its-surface:tut}

\begin{tutorial}
\textbf{Goal.} Build a surface from quotes in total variance, check it for \omterm{def:dv:implied-volatility-and-its-surface:static}{static arbitrage},
break it with a bad quote, and read the \omterm{def:dv:implied-volatility-and-its-surface:density}{risk-neutral density} off it. \textbf{End state:}
\cref{fig:dv:implied-volatility-and-its-surface:smiles,fig:dv:implied-volatility-and-its-surface:density}
and the numbers of \cref{ex:dv:implied-volatility-and-its-surface:crash}.
\begin{tutsteps}
\item \textbf{The checks}: the density factor $g$ on every slice, and $w$ across slices.
\omcode{code/firm/volsurface/firm_volsurface.py}{102}{120}{Butterfly and calendar checks in total variance.}\label{lst:dv:implied-volatility-and-its-surface:check}
\item \textbf{The crash premium}: the probability below a level is the strike derivative of
the put, so the skew enters through the derivative of the volatility.
\omcode{code/derivatives/07-implied-volatility-and-its-surface/python/dv_surface.py}{80}{95}{Risk-neutral probability of a fall, with and without the skew.}\label{lst:dv:implied-volatility-and-its-surface:crash}
\item \textbf{Run} \texttt{dv\_surface.check(index\_surface())} (no violation),
\texttt{check(index\_surface((1, -0.1, 0.03)))} (six), \texttt{crash\_probability()} and
\texttt{fig\_surface.py}.
\end{tutsteps}
\textbf{What to change next.} Interpolate each smile linearly in $k$ instead of by a cubic
spline and count the butterfly violations the kinks create; then build the surface on a strike
grid rather than a moneyness grid and look for calendar violations near ex-dividend dates.
\end{tutorial}

\section{Build: the volatility surface}\label{bld:dv:implied-volatility-and-its-surface:volsurface}

\begin{build}
\bfield{Purpose} The miniature firm's single object for implied volatility: fits (chapter 8),
models (chapters 9--13) and the pricing library (chapter 28) read from it, and nothing downstream
of it may see an arbitrage.
\bfield{Interface} \texttt{Surface(asof, expiries, ks, w, forwards)} with
\texttt{implied\_vol(strike, expiry)}, \texttt{total\_variance(k, t)}, \texttt{forward(t)};
\texttt{from\_vols(asof, expiries, forwards, ks, vols)}; \texttt{arbitrage\_report(surface)};
\texttt{density\_factor}; \texttt{risk\_neutral\_density(strikes, calls, df)}.
\bfield{Rules} Interpolate total variance: a not-a-knot cubic spline in $k$ within an expiry
(a natural spline bends a convex smile the wrong way at its ends), linear in $t$ at fixed $k$
between expiries, flat volatility before the first and after the last; extrapolate linearly in
$k$ and report the wings as a risk (chapter 8).
\bfield{Acceptance tests} \texttt{code/firm/volsurface/tests/}: the surface returns its input
volatilities at the nodes; a surface built from a flat volatility has a density factor of one
at $k=0$ and no violation; a bumped quote creates butterfly violations and a lowered long
expiry a calendar violation; Breeden--Litzenberger recovers the lognormal density.
\bfield{Stretch} Replace the node grid by the SVI slices of chapter 8, and add a sticky-delta
mode that re-centres the surface on a new forward.
\end{build}

\begin{omsources}
\item E.~Derman, \emph{The Volatility Smile}, lecture notes, Columbia University (2008),
lecture 1; and ``Regimes of volatility'', \emph{Risk} (April 1999).
\item M.~Rubinstein, ``Implied binomial trees'', \emph{Journal of Finance} 49 (1994) 771--818.
\item D.~T.~Breeden and R.~H.~Litzenberger, ``Prices of state-contingent claims implicit in
option prices'', \emph{Journal of Business} 51 (1978) 621--651.
\item P.~Carr and D.~B.~Madan, ``A note on sufficient conditions for no arbitrage'',
\emph{Finance Research Letters} 2 (2005) 125--130.
\item G.~Bakshi, N.~Kapadia and D.~Madan, ``Stock return characteristics, skew laws, and the
differential pricing of individual equity options'', \emph{Review of Financial Studies} 16
(2003) 101--143.
\item Y.-F.~Chen and X.~Mu, ``Asymmetric volatility in commodity markets'', \emph{Journal of
Commodity Markets} 22 (2021).
\item J.~Gatheral, \emph{The Volatility Surface}, Wiley, 2006.
\end{omsources}

\section{Exercises}

\begin{exercise}[$\star$]\label{exo:dv:implied-volatility-and-its-surface:1}
The three-month forward is 100.37. Give the \omterm{def:dv:implied-volatility-and-its-surface:surface}{log-moneyness} of the strike 90 and its total
variance at 20.47 volatility.
\end{exercise}

\begin{exercise}[$\star$]\label{exo:dv:implied-volatility-and-its-surface:2}
On the chapter's surface the three-month 90 and 110 strikes are at 20.47 and 14.19
volatility. Give the 90--110 skew and explain why the same skew in strike terms would be
smaller at one year.
\end{exercise}

\begin{exercise}[$\star$]\label{exo:dv:implied-volatility-and-its-surface:3}
At the same forward-moneyness $k=0$, the three-month option is quoted at 20\% and the
six-month at 13\%. Is there a \omterm{def:dv:implied-volatility-and-its-surface:static}{calendar arbitrage}? What would you trade?
\end{exercise}

\begin{exercise}[$\star\star$]\label{exo:dv:implied-volatility-and-its-surface:4}
Price the three-month at-the-money digital call on the chapter's surface with and without
the skew correction, given a \omterm{def:dv:greeks-and-the-hedging-pnl:greeks}{vega} of 19.77 and a volatility slope $\partial\sigma/\partial K$
of $-0.0031$ per unit of strike.
\end{exercise}

\begin{exercise}[$\star\star$]\label{exo:dv:implied-volatility-and-its-surface:5}
Give the three-month call's delta at the 100 strike under \omterm{def:dv:implied-volatility-and-its-surface:sticky}{sticky strike} and under \omterm{def:dv:implied-volatility-and-its-surface:sticky}{sticky
delta}, and say how many shares per 100 calls the choice changes.
\end{exercise}

\begin{exercise}[$\star\star$]\label{exo:dv:implied-volatility-and-its-surface:6}
Why is an equity index's skew steeper than the skews of its members?
\end{exercise}

\begin{exercise}[$\star\star\star$]\label{exo:dv:implied-volatility-and-its-surface:7}
\emph{Coding.} With \texttt{index\_surface} and \texttt{check}, add a quote error of $+3$
points at three months and $k=-0.1$: how many butterfly violations, and how negative does
$g$ get? Then find the smallest downward error on the six-month at-the-money quote that
creates a calendar violation.
\end{exercise}

\begin{exercise}[$\star\star\star$]\label{exo:dv:implied-volatility-and-its-surface:8}
\emph{Find the flaw.} ``The three-month 80 put is at 27 volatility against 16 at the money:
the market expects the index to be twice as volatile if it falls to 80.''
\end{exercise}

\section{Problem: The Crash Premium}

\begin{problem}[{Weekend problem --- what the skew says about a 20\% fall}]\label{pb:dv:implied-volatility-and-its-surface:1}
A pension fund asks what the index options say about a fall of more than 20\% over the next
three months, and what protection against it costs. Use the chapter's surface: spot 100,
rate 3\%, dividend yield 1.5\%, three months to expiry.

\medskip\noindent\textbf{Part I --- The inputs.}
\begin{enumerate}
\item Give the three-month forward and discount factor.
\item Give the at-the-money volatility and the volatility at the 80 strike.
\item Give the \omterm{def:dv:implied-volatility-and-its-surface:surface}{log-moneyness} of the 80 strike.
\item Price the 80 put on the skewed surface.
\item Price it on a flat surface at the at-the-money volatility.
\end{enumerate}

\medskip\noindent\textbf{Part II --- Probabilities.}
\begin{enumerate}[resume]
\item Give the risk-neutral probability of ending below 80 on the flat surface.
\item Give it on the skewed surface.
\item Why is it not $\Phi(-d_2)$ at the 80 strike's own volatility?
\item How is the probability related to a digital put, and to a put spread?
\item Sketch where the skew moves probability, using the density.
\end{enumerate}

\medskip\noindent\textbf{Part III --- Interpreting.}
\begin{enumerate}[resume]
\item Is 1.89\% a forecast of the frequency of such falls? Why not?
\item What else is in the price of the 80 put?
\item How would a 5\% fall of the index change the 80 put under \omterm{def:dv:implied-volatility-and-its-surface:sticky}{sticky strike} and \omterm{def:dv:implied-volatility-and-its-surface:sticky}{sticky delta}?
\item Why did the skew appear after 1987?
\item Why might the skew be steeper at one month than at one year?
\end{enumerate}

\medskip\noindent\textbf{Part IV --- Judgement.}
\begin{enumerate}[resume]
\item What would you tell the fund about the cost of this protection?
\item What is the cheaper way to buy protection against a 20\% fall, and what does it give up?
\item What should a risk system use for this probability: the flat or the skewed surface?
\item State the \emph{named result}: the risk-neutral probability of a fall below 80 in three
months on the flat and on the skewed surface, and their ratio.
\item In one sentence: what does the skew price?
\end{enumerate}
\end{problem}

\section{Interview questions}

\begin{interviewq}[$\star$ \iqroles{trader, researcher}]\label{iq:dv:implied-volatility-and-its-surface:1}
Why do equity index options have a skew, and why did it appear only after 1987?
\end{interviewq}

\begin{interviewq}[$\star$ \iqroles{researcher}]\label{iq:dv:implied-volatility-and-its-surface:2}
How do you get the \omterm{def:dv:implied-volatility-and-its-surface:density}{risk-neutral density} from option prices?
\end{interviewq}

\begin{interviewq}[$\star\star$ \iqroles{trader}]\label{iq:dv:implied-volatility-and-its-surface:3}
The index falls 5\%. Under \omterm{def:dv:implied-volatility-and-its-surface:sticky}{sticky strike} and under \omterm{def:dv:implied-volatility-and-its-surface:sticky}{sticky delta}, what happens to the
at-the-money volatility and to your delta?
\end{interviewq}

\begin{interviewq}[$\star\star$ \iqroles{researcher, developer}]\label{iq:dv:implied-volatility-and-its-surface:4}
Give the static no-arbitrage conditions on an implied-volatility surface, and how you check
them.
\end{interviewq}

\begin{interviewq}[$\star\star$ \iqroles{risk}]\label{iq:dv:implied-volatility-and-its-surface:5}
A risk report computes the probability of a 20\% fall from the at-the-money volatility. What
is wrong, and by how much can it be off?
\end{interviewq}

\begin{interviewq}[$\star\star\star$ \iqroles{developer, researcher}]\label{iq:dv:implied-volatility-and-its-surface:6}
How would you interpolate a surface between quoted strikes and expiries without creating
arbitrage?
\end{interviewq}
